% concatenated from main.tex
\documentclass[11pt]{amsart}

\usepackage{amsmath,amssymb,amsthm,mathtools}
\usepackage{a4wide}
\usepackage{tikz-cd}
\usepackage[colorlinks,citecolor=blue,urlcolor=blue]{hyperref}
\usepackage{cleveref}

\renewcommand{\leq}{\leqslant}
\renewcommand{\geq}{\geqslant}
\renewcommand{\le}{\leqslant}
\renewcommand{\ge}{\geqslant}

\numberwithin{equation}{section}

% Theorem environments
\newtheorem{theorem}{Theorem}[section]
\newtheorem{lemma}[theorem]{Lemma}
\newtheorem{proposition}[theorem]{Proposition}
\newtheorem{corollary}[theorem]{Corollary}

\theoremstyle{definition}
\newtheorem{definition}[theorem]{Definition}
\newtheorem{example}[theorem]{Example}
\newtheorem{question}[theorem]{Question}

\theoremstyle{remark}
\newtheorem{remark}[theorem]{Remark}

% Main theorems labelled A, B, C, ...
\theoremstyle{plain}
\newcounter{maintheorem}
\renewcommand*{\themaintheorem}{\Alph{maintheorem}}
\newtheorem{mainth}[maintheorem]{Theorem}
\crefname{maintheorem}{Theorem}{Theorems}

% Operators
\DeclareMathOperator{\Hom}{Hom}
\DeclareMathOperator{\supp}{supp}
\DeclareMathOperator{\id}{id}

\begin{document}

\title[Semiprojective Banach lattices]{Semiprojective Banach lattices}

\author[T.~Kania]{Tomasz Kania}
\address{Institute of Mathematics, Czech Academy of Sciences,
\v{Z}itn\'{a} 25, 115~67 Prague 1, Czech Republic
\& Institute of Mathematics, Jagiellonian University,
{\L}ojasiewicza 6, 30-348 Krak\'ow, Poland}
\email{tomasz.marcin.kania@gmail.com}
\thanks{Supported by IM CAS (RVO 67985840).}

\author[M.~Niwi\'nski]{Mariusz Niwi\'nski}
\address{Doctoral School of Exact and Natural Sciences,
Faculty of Mathematics and Computer Science,
Institute of Mathematics, Jagiellonian University,
{\L}ojasiewicza 6, 30-348 Krak\'ow, Poland}
\email{mariusz.niwinski@doctoral.uj.edu.pl}

\date{\today}

\subjclass[2020]{46B42, 46E15, 54C15, 18A30}
\keywords{Banach lattice, semiprojectivity, inductive limit ideal, absolute neighbourhood retract,
$C(X)$-space, finitely presentable object, $\ell$-open}

\begin{abstract}
We introduce a norm-controlled notion of semiprojectivity for Banach lattices,
requiring liftability of contractive lattice homomorphisms through inductive limits of closed ideals
with arbitrarily small loss of norm control.
Our main result establishes that, for a~compact metric space~$X$,
the Banach lattice $C(X)$ is semiprojective if and only if $X$ is an absolute neighbourhood retract.

Notably, this characterisation is strictly more permissive than its $C^*$-algebraic counterpart:
by a theorem of S\o rensen and Thiel,
$C(X)$ is semiprojective in the category of $C^*$-algebras and $*$-homo\-morphisms
if and only if $X$ is an ANR of dimension at most one.
The dimensional obstruction disappears in the Banach-lattice setting because lattice homomorphisms
between $C(K)$-spaces are automatically weighted composition operators,
and therefore no commutation relations need to be lifted.
We also show that uncountable $\ell_1$-sums of $1^+$-projective Banach lattices
with topological order units are semiprojective but need not be $1^+$-projective,
establishing that the two notions are genuinely distinct.
On the negative side, we prove that $\ell_p$ \((1<p<\infty)\) and \(L_p([0,1])\) \((1\le p<\infty)\) as well as Orlicz spaces
are not semiprojective.
\end{abstract}

\maketitle

\section{Introduction}\label{sec:intro}

Projective objects encode extension and lifting phenomena in a single categorical property:
every morphism into a quotient factors through the quotient map.
In the Banach-space world the strictly isometric version of this requirement
is overly rigid, and it is by now standard to relax it to an \emph{approximate} lifting condition
in which the norm of the factoring map is allowed to exceed the original norm
by an arbitrarily small factor $1+\varepsilon$.
This philosophy underlies the notion of a \emph{$1^+$-projective} Banach lattice introduced by
de~Pagter and Wickstead~\cite{dpw}, who showed, among other things,
that $\ell_1$ is $1^+$-projective.
More recently, Avil\'es, Mart\'{\i}nez-Cervantes and Rodr\'{\i}guez Abell\'an~\cite{avilaes2020}
characterised the $1^+$-projective spaces of the form $C(K)$:
\emph{$C(K)$ is $1^+$-projective if and only if $K$ is an absolute neighbourhood retract (ANR)}.

In many situations, however, one need not lift through arbitrary quotients
but only through those arising from increasing unions of ideals---that is, through
\emph{inductive limits}.
The resulting \emph{semiprojective} lifting property is considerably weaker than full projectivity,
yet retains substantial force.
Semiprojectivity and its variants play a central r\^ole in the theory of
$C^*$-algebras, where the notion was introduced by Blackadar~\cite{blackadar85}
and systematically developed in the work of Loring~\cite{loring} and others.
In the commutative $C^*$-setting,
semiprojectivity of $C(X)$ is intimately connected with $X$ being an ANR,
but an additional \emph{dimensional} constraint intervenes:
S\o rensen and Thiel~\cite{sorensen-thiel} proved that
\emph{$C(X)$ is semiprojective in the category of $C^*$-algebras
if and only if $X$ is an ANR with $\dim X\le 1$},
confirming a conjecture of Blackadar~\cite[p.~24]{blackadar-survey}.
The dimensional obstruction arises because commutation relations---which constrain $*$-homomorphisms
from $C(X)$ into a noncommutative ambient algebra---are notoriously hard to lift;
for instance, $C(\mathbb T^2)$ is not semiprojective in the $C^*$-category,
as can be shown using the Voiculescu matrices (see, \emph{e.g.},
\cite[p.~23]{blackadar-survey}).

In a recent survey, Thiel~\cite{thiel-categories} studies semiprojectivity and its
dual notion, semiinjectivity, as general categorical concepts.
In the category of compact metric spaces and continuous maps,
(semi)injective objects are precisely the absolute (neighbourhood) retracts;
thus semiinjectivity of~$X$ in that category corresponds, via Gelfand duality,
to semiprojectivity of~$C(X)$ among commutative $C^*$-algebras.
In the category of groups, the picture is strikingly different:
a group is semiprojective if and only if it is a retract of a free product
of a finitely presented group and a free group, while the trivial group is
the only semiinjective group~\cite[Propositions~3.7 and~3.9]{thiel-categories}.
Thiel's framework highlights that the behaviour of semiprojectivity is
strongly category-dependent, making it natural to investigate the concept in the
Banach-lattice category, where the morphisms---lattice homomorphisms with
controlled norms---differ fundamentally from both $*$-homomorphisms and group
homomorphisms.

The present note records the analogous phenomenon in the Banach-lattice category.
Our main results read as follows.

\begin{mainth}\label{thm:A}
Let $X$ be a compact metric space.
Then the Banach lattice $C(X)$ is semiprojective
if and only if $X$ is an absolute neighbourhood retract.
\end{mainth}

Here semiprojectivity is understood in the sense of Definition~\ref{def:semi} below.
The ``if'' direction is immediate: if $X$ is an ANR, then $C(X)$ is $1^+$-projective
by~\cite[Theorem~1.4]{avilaes2020}, and $1^+$-projectivity evidently implies semiprojectivity.
The content of the theorem lies in the converse:
if $X$ is not an ANR, then $C(X)$ is not semiprojective.
The proof, given in Section~\ref{sec:CX}, exploits the representation of lattice
homomorphisms between $C(K)$-spaces as weighted composition operators
(Lemma~\ref{lem:weighted-composition}).
A hypothetical semiprojective lifting produces, via this representation,
a continuous retraction from a neighbourhood of $X$ onto $X$,
contradicting the assumption that $X$ is not an ANR.

\begin{mainth}\label{thm:B}
Let\/ $\Gamma$ be a set and let\/ $(P_\gamma)_{\gamma\in\Gamma}$ be a family of
$1^+$-projective Banach lattices, each admitting a topological order unit.
Then the Banach lattice
\[
P:=\Bigl(\bigoplus_{\gamma\in\Gamma} P_\gamma\Bigr)_{\ell_1(\Gamma)}
\]
is semiprojective. If\/ $\Gamma$ is countable, then $P$ is $1^+$-projective.
\end{mainth}

Recall that an element $e\ge 0$ in a Banach lattice~$E$ is a
\emph{topological order unit} if the closed ideal it generates equals all of~$E$,
that is, if every $x\in E$ can be approximated by elements $y$ satisfying $|y|\le c\,e$
for some constant~$c>0$.
Every $C(K)$-space admits a topological order unit, namely the constant function
$\mathbf{1}_K$, since $|f|\le \|f\|_\infty\,\mathbf{1}_K$ for all $f\in C(K)$.

\begin{corollary}\label{cor:l1-Gamma-intro}
For every set\/ $\Gamma$, the Banach lattice $\ell_1(\Gamma)$ is semiprojective.
It is $1^+$-projective if and only if\/ $\Gamma$ is countable.
\end{corollary}

In particular, Corollary~\ref{cor:l1-Gamma-intro} shows that semiprojectivity and
$1^+$-projectivity are genuinely distinct notions for Banach lattices.
The key observation is that any lattice homomorphism from the $\ell_1$-sum
into a separable target annihilates all but countably many coordinates
(the images of the coordinate order units form a disjoint family,
and disjoint families in separable Banach lattices are countable),
so the lifting problem reduces to a countable $\ell_1$-sum of $1^+$-projectives.

On the other hand, we show that $\ell_p$ is \emph{not} semiprojective for $1<p<\infty$
(Corollary~\ref{cor:ell-p-not-semi}),
using a path-space construction that converts semiprojectivity into an exact lifting
property and then applying a growth estimate of
Avil\'es, Mart\'{\i}nez-Cervantes and Rodr\'{\i}guez Abell\'an~\cite{avilaes-c0}.

\subsection*{Comparison with the $C^*$-algebraic setting}
It is instructive to contrast Theorem~\ref{thm:A} with the S\o rensen--Thiel
result~\cite{sorensen-thiel}.
In the $C^*$-category, semiprojectivity of $C(X)$ requires $\dim X\le 1$
\emph{in addition to} $X$ being an ANR;
in the Banach-lattice category the dimension plays no r\^ole.
The reason is structural:
a lattice homomorphism $T\colon C(X)\to C(Y)$ is automatically a weighted composition operator
(Lemma~\ref{lem:weighted-composition}), determined by a continuous weight and a continuous map $Y\supseteq U\to X$.
In particular, $T$ is completely determined by its values on one-dimensional fibres,
and no commutation relations need to be preserved.
By contrast, a $*$-homomorphism $C(X)\to B$ for a general $C^*$-algebra $B$
must respect commutativity of $C(X)$, and lifting such commutation constraints
through inductive limits fails for $\dim X\ge 2$, as the Voiculescu-matrix obstruction
demonstrates.

This makes the Banach-lattice category a natural setting in which
the \emph{purely topological} characterisation
``semiprojective $\Longleftrightarrow$ ANR'' holds without dimensional qualifications,
mirroring the classical shape-theoretic equivalence more faithfully
than the $C^*$-category does.
The following table summarises the comparison (all spaces are compact and metrisable).

\medskip
\begin{center}
\renewcommand{\arraystretch}{1.3}
\begin{tabular}{lcc}
\hline
\textbf{Property of $C(X)$} & \textbf{Banach lattices} & \textbf{$C^*$-algebras}\\
\hline
$1^+$-projective / projective & $X$ is an ANR & $X$ is an AR, $\dim X\le 1$ \\
Semiprojective & $X$ is an ANR & $X$ is an ANR, $\dim X\le 1$ \\
\hline
\end{tabular}
\end{center}
In the Banach-lattice column, the relevant notion is \(1^+\)-projectivity
(which is called simply \emph{projective} in \cite{dpw});
in the \(C^*\)-column, \emph{projective} is used in the usual exact sense.

For commutative $C^*$-algebras, projectivity of $C(X)$
was characterised by Chigogidze and Dranishnikov~\cite{chigogidze-dranishnikov},
and semiprojectivity by S\o rensen and Thiel~\cite{sorensen-thiel}.

\medskip

\noindent\textbf{Notation and conventions.}
All Banach lattices are real.
A~\emph{lattice homomorphism} between Banach lattices
is a bounded linear map $T$ satisfying $T|x|=|Tx|$ for every~$x$.
\emph{Closed ideals} are understood in the lattice sense:
a closed subspace $J$ of a Banach lattice $B$ is a (lattice) ideal
if $|b|\le |a|$ and $a\in J$ imply $b\in J$.
For a compact Hausdorff space~$K$ we write $C(K)$
for the Banach lattice of continuous real-valued functions on~$K$
equipped with the supremum norm and the pointwise ordering.
If $F\subseteq K$ is closed, we set $I_F:=\{f\in C(K): f|_F=0\}$.

\section{Semiprojectivity and norm control}\label{sec:semi}

We begin by recording the definitions that will be used throughout.

\begin{definition}[Inductive limit ideal]\label{def:ili}
Let $B$ be a Banach lattice.
A closed ideal $I\subseteq B$ is an \emph{inductive limit ideal} if there exists
an increasing sequence of closed ideals $I_1\subseteq I_2\subseteq \cdots\subseteq B$ such that
\[
I=\overline{\bigcup_{n=1}^\infty I_n}.
\]
We write $Q_n\colon B\to B/I_n$ for the quotient map and
$\pi_n\colon B/I_n\to B/I$ for the canonical surjection induced by the
inclusion $I_n\subseteq I$.
\end{definition}

\begin{definition}[Semiprojectivity]\label{def:semi}
A Banach lattice $A$ is \emph{semiprojective} if for every \emph{separable} Banach lattice~$B$,
every inductive limit ideal $I=\overline{\bigcup_n I_n}\subseteq B$,
and every contractive lattice homomorphism $\phi\colon A\to B/I$,
there exist $n\in\mathbb N$ and, for every $\varepsilon>0$,
a lattice homomorphism $\psi\colon A\to B/I_n$ satisfying
\[
\|\psi\|\le 1+\varepsilon
\qquad\text{and}\qquad
\pi_n\circ \psi=\phi.
\]
In other words, $\phi$ eventually lifts through the tower of quotients with arbitrarily
small loss of norm control, as depicted in the following diagram:
\[
\begin{tikzcd}[row sep=1.6em, column sep=3em]
B \arrow[d, two heads, "Q_1"'] \\
B/I_1 \arrow[d, two heads] \\
\vdots \arrow[d, two heads] \\
B/I_n \arrow[d, two heads, "\pi_n"'] & A \arrow[l, dashed, "\psi"'] \arrow[dl, "\phi"] \\
B/I
\end{tikzcd}
\qquad\quad \pi_n\circ\psi=\phi,\quad \|\psi\|\le 1+\varepsilon.
\]
Here the solid vertical arrows are the canonical surjections,
$\phi$ is a given contractive lattice homomorphism, and the dashed arrow $\psi$ is
the required approximate lift satisfying $\pi_n\circ \psi=\phi$.
Crucially, the index $n$ must be chosen \emph{independently} of $\varepsilon>0$.
\end{definition}

\begin{remark}\label{rem:countable-exhaustion}
If $B$ is separable and $I\subseteq B$ is a closed ideal, then $I$ admits an
increasing exhaustion by closed ideals. Indeed, choose a dense sequence
$(x_n)_{n\ge 1}$ in $I$, and let $I_n$ be the closed ideal generated by
$\{x_1,\dots,x_n\}$. Then
\[
I=\overline{\bigcup_{n=1}^\infty I_n}.
\]
Thus, once the ambient Banach lattice is separable, the adjective
``inductive limit ideal'' does not restrict the class of ideals under
consideration; what remains essential in Definition~\ref{def:semi} is the
chosen approximating sequence of ideals.
\end{remark}

\begin{remark}\label{rem:finitely-presentable}
In the theory of locally presentable categories
(see Ad\'amek and Rosick\'y~\cite[Definition~1.1]{adamek-rosicky}),
an object~$K$ of a category~$\mathcal C$ is \emph{finitely presentable}
if the covariant hom-functor
$\Hom_{\mathcal C}(K,{-})$ preserves directed colimits.
For algebraic categories (varieties of algebras) this amounts to
$K$ being presented by finitely many generators
and finitely many relations;
in particular, the quotient of a free algebra by a finitely generated
ideal possesses the relevant lifting property with respect to directed unions.

As observed by L.~Positselski, Definition~\ref{def:semi} can be viewed as a
norm-controlled, approximate analogue of finite presentability.
The r\^ole of the directed colimit is played by
the inductive limit ideal $I=\overline{\bigcup I_n}$,
and the factoring morphism is required to exist only up to
an arbitrarily small norm defect~$\varepsilon$.
The restriction to \emph{sequential} inductive limits (rather than arbitrary directed families)
is a convenient separable variant, reflecting that many norm-approximation and lifting arguments
in Banach-space settings are naturally formulated along sequences.
It would be interesting to develop this analogy further within a systematic framework for
``approximate locally presentable categories'' of Banach lattices;
we leave this for future investigation.
\end{remark}

\begin{remark}[A categorical perspective]\label{rem:cat}
For each $\lambda\ge 1$ one can form a category $\mathbf{BL}_\lambda$
whose objects are Banach lattices and whose morphisms are lattice homomorphisms
of norm at most~$\lambda$.
The family $(\mathbf{BL}_\lambda)_{\lambda\ge 1}$ is \emph{multiplicatively graded}:
the composition of a $\lambda$-morphism with a $\mu$-morphism is a $(\lambda\mu)$-morphism.
The ubiquitous factor $1+\varepsilon$ in norm-controlled lifting problems
thus reflects the freedom to work in any category $\mathbf{BL}_\lambda$ with $\lambda>1$.

To eliminate the quantifier over $\varepsilon$ one can pass to an auxiliary
\emph{approximately contractive} category $\mathbf{BL}_{1^+}$.
Its objects are Banach lattices, and a morphism $E\to F$ is a sequence
$(T_k)_{k\ge 1}$ of lattice homomorphisms $T_k\colon E\to F$ such that
$\|T_k\|\le 1+\varepsilon_k$ for some real sequence $\varepsilon_k\to 0$.
Composition is defined termwise:
\[
(S_k)_{k\ge 1}\circ (T_k)_{k\ge 1}:=(S_k\circ T_k)_{k\ge 1}.
\]
This is well-defined since, if $\|T_k\|\le 1+\varepsilon_k$ and
$\|S_k\|\le 1+\delta_k$ with $\varepsilon_k,\delta_k\to 0$, then
\[
\|S_k\circ T_k\|\le (1+\delta_k)(1+\varepsilon_k)=1+\gamma_k,
\qquad
\gamma_k:=\varepsilon_k+\delta_k+\varepsilon_k\delta_k\to 0.
\]
Every contractive lattice homomorphism~$T$ embeds into $\mathbf{BL}_{1^+}$
via the constant sequence $(T,T,\dots)$.
In this language, $A$ is semiprojective precisely when every contractive
$\phi\colon A\to B/I$ lifts, for some~$n$,
to a morphism $(\psi_k)\colon A\to B/I_n$ in $\mathbf{BL}_{1^+}$
satisfying $\pi_n\circ \psi_k=\phi$ for all~$k$.
\end{remark}

\begin{proposition}[Retracts preserve semiprojectivity]\label{prop:retract}
Let $A$ and $P$ be Banach lattices and suppose that there exist \emph{contractive}
lattice homomorphisms $i\colon A\to P$ and $r\colon P\to A$ with
$r\circ i=\id_A$.
If $P$ is semiprojective, then $A$ is semiprojective.
\end{proposition}

\begin{proof}
Let $B$ be a separable Banach lattice and let
$I=\overline{\bigcup_{n=1}^\infty I_n}\subseteq B$ be an inductive limit ideal.
Let $\phi\colon A\to B/I$ be a contractive lattice homomorphism.
Then $\phi\circ r\colon P\to B/I$ is a contractive lattice homomorphism.
Since $P$ is semiprojective, there exists $n\in\mathbb N$ such that for every
$\varepsilon>0$ there is a lattice homomorphism
$\widetilde\psi_\varepsilon\colon P\to B/I_n$ satisfying
$\|\widetilde\psi_\varepsilon\|\le 1+\varepsilon$ and
$\pi_n\circ \widetilde\psi_\varepsilon=\phi\circ r$.
Set $\psi_\varepsilon:=\widetilde\psi_\varepsilon\circ i\colon A\to B/I_n$.
Then $\psi_\varepsilon$ is a lattice homomorphism,
$\|\psi_\varepsilon\|\le 1+\varepsilon$, and
\[
\pi_n\circ \psi_\varepsilon
=\pi_n\circ \widetilde\psi_\varepsilon\circ i
=\phi\circ r\circ i
=\phi.
\]
Thus $A$ is semiprojective.
\end{proof}

\section{A topological obstruction for $C(X)$}\label{sec:CX}

In this section we prove the non-trivial direction of Theorem~\ref{thm:A}:
if $X$ is a compact metric space which is not an ANR,
then $C(X)$ is not semiprojective.

\subsection{Ideals in $C(K)$}

If $F$ is a closed subset of a compact Hausdorff space~$K$,
the restriction map $\rho_F\colon C(K)\to C(F)$ is a surjective
contractive lattice homomorphism with kernel~$I_F$,
hence it induces an isometric lattice isomorphism $C(K)/I_F\cong C(F)$.

\begin{lemma}\label{lem:limit-ideal}
Let $Y$ be a compact metric space, let $X\subseteq Y$ be closed,
and fix a decreasing sequence of compact neighbourhoods
$Y_1\supseteq Y_2\supseteq\cdots\supseteq X$
with $\bigcap_{n} Y_n=X$.
Set $I:=I_X$ and $I_n:=I_{Y_n}$ in $C(Y)$.
Then $I=\overline{\bigcup_{n=1}^\infty I_n}$.
\end{lemma}

\begin{proof}
Since $I_n\subseteq I$ for all~$n$, the inclusion
$\overline{\bigcup_n I_n}\subseteq I$ is clear.
For the reverse, let $f\in I$ and $\varepsilon>0$.
As $f$ vanishes on~$X$ and is continuous,
there exists an open set $U\supseteq X$ in $Y$ such that $|f(y)|<\varepsilon$ for all $y\in U$.
Choose $n$ with $Y_n\subseteq U$.
Such an $n$ exists since $(Y_n)$ is a decreasing sequence of compact sets with
$\bigcap_n Y_n=X\subseteq U$:
otherwise one could pick $y_n\in Y_n\setminus U$ and extract a convergent subsequence
with limit in $\bigcap_n Y_n=X$, contradicting $X\subseteq U$.
By Urysohn's lemma (applied in the normal space~$Y$) there is a
function $u\in C(Y)$ with $0\le u\le 1$, $u|_{Y_n}=0$ and $u|_{Y\setminus U}=1$.
Set $g:=uf$.
Then $g|_{Y_n}=0$, so $g\in I_n$, and
\[
\|f-g\|_\infty=\sup_{y\in Y}\bigl|(1-u(y))f(y)\bigr|\le \varepsilon,
\]
since $1-u$ is supported on~$U$, where $|f|<\varepsilon$.
Thus $f\in\overline{\bigcup_n I_n}$.
\end{proof}

\subsection{Lattice homomorphisms between $C(K)$-spaces}

The following lemma is standard; we include a short proof for the reader's convenience.

\begin{lemma}\label{lem:eval}
Let $X$ be a compact Hausdorff space and let $\varphi\colon C(X)\to \mathbb R$ be a non-zero lattice
homomorphism.
Then there exist unique $x\in X$ and $c>0$ such that $\varphi(f)=c\,f(x)$ for all $f\in C(X)$.
\end{lemma}

\begin{proof}
Set $c:=\varphi(\mathbf{1})>0$ (where $\mathbf{1}$ denotes the constant function~$1$)
and $\psi:=c^{-1}\varphi$.
Then $\psi$ is a unital positive linear functional of norm~$1$, so by the
Riesz--Markov representation theorem there is a Borel probability measure $\mu$ on~$X$
with $\psi(f)=\int_X f\,d\mu$ for all $f\in C(X)$.

If $\mu$ were not a point mass, then $\supp(\mu)$ would contain two distinct points
$x_1\ne x_2$.
Choose disjoint open neighbourhoods $U,V\subseteq X$ of $x_1$ and $x_2$.
Since $x_1,x_2\in\supp(\mu)$ we have $\mu(U)>0$ and $\mu(V)>0$.
By inner regularity of $\mu$ there exist compact sets $K\subseteq U$ and $L\subseteq V$
with $\mu(K)>0$ and $\mu(L)>0$.
By Urysohn's lemma, choose $f,g\in C(X)_+$ such that
$f\le \mathbf{1}$, $g\le \mathbf{1}$,
$f|_{K}\equiv 1$, $g|_{L}\equiv 1$,
and $\supp(f)\subseteq U$, $\supp(g)\subseteq V$.
Then $f\wedge g=0$, whence
\[
0=\psi(f\wedge g)=\psi(f)\wedge \psi(g)
=\min\bigl\{\psi(f),\,\psi(g)\bigr\}.
\]
On the other hand,
\[
\psi(f)=\int_X f\,d\mu\ge \mu(K)>0,
\qquad
\psi(g)=\int_X g\,d\mu\ge \mu(L)>0,
\]
a contradiction.
Therefore $\mu=\delta_x$ for a unique $x\in X$,
so $\varphi(f)=c\,f(x)$ for all $f\in C(X)$.
\end{proof}

The next lemma is also well known but we present it for the sake of completeness.

\begin{lemma}\label{lem:weighted-composition}
Let $X$ and $Y$ be compact Hausdorff spaces and let $T\colon C(X)\to C(Y)$ be a~lattice homomorphism.
Set $h:=T\mathbf{1}\in C(Y)_+$ and $U:=\{y\in Y: h(y)>0\}$.
Then there exists a~unique continuous map $r\colon U\to X$ such that
\begin{equation}\label{eq:wc}
(Tf)(y)=
\begin{cases}
h(y)\,f\bigl(r(y)\bigr) & \text{if } y\in U,\\[2pt]
0 & \text{if } y\in Y\setminus U,
\end{cases}
\end{equation}
for all $f\in C(X)$.
\end{lemma}

\begin{proof}
For each $y\in Y$ define $\varphi_y\colon C(X)\to\mathbb R$ by $\varphi_y(f):=(Tf)(y)$.
Then $\varphi_y$ is a lattice homomorphism.
If $h(y)=0$, then for every $f\in C(X)_+$ we have
$0\le \varphi_y(f)\le \|f\|_\infty\,\varphi_y(\mathbf{1})=0$,
so $\varphi_y(f)=0$.
By linearity it follows that $\varphi_y=0$, hence $(Tf)(y)=0$ for all~$f$.
If $y\in U$, then $\varphi_y\ne 0$ and Lemma~\ref{lem:eval} provides
unique $c(y)>0$ and $r(y)\in X$ with $\varphi_y(f)=c(y)\,f(r(y))$.
Evaluating at $f=\mathbf{1}$ gives $c(y)=h(y)$, establishing~\eqref{eq:wc}.
Uniqueness of $r$ follows from Lemma~\ref{lem:eval}.

It remains to verify that $r\colon U\to X$ is continuous.
Suppose $(y_\alpha)$ is a net in $U$ converging to some $y\in U$
with $r(y_\alpha)\to x\in X$ along a subnet (compactness of~$X$ guarantees such subnets exist).
For any $f\in C(X)$, continuity of $Tf$ and $h$ yields
\[
f(x)
=\lim_\alpha f\bigl(r(y_\alpha)\bigr)
=\lim_\alpha \frac{(Tf)(y_\alpha)}{h(y_\alpha)}
=\frac{(Tf)(y)}{h(y)}
=f\bigl(r(y)\bigr).
\]
Since $f$ was arbitrary, $x=r(y)$.
Thus every convergent subnet of $(r(y_\alpha))$ has limit $r(y)$.
In particular, $(r(y_\alpha))$ has the unique cluster point $r(y)$.
Since $X$ is compact Hausdorff, this implies $r(y_\alpha)\to r(y)$,
and $r$ is continuous.
\end{proof}

\subsection{Proof of Theorem~\ref{thm:A}}

\begin{proof}[Proof of Theorem~\ref{thm:A}]
As noted in the introduction, if $X$ is an ANR then $C(X)$ is $1^+$-projective
by~\cite[Theorem~1.4]{avilaes2020}, hence semiprojective.
We prove the converse.

Suppose that $X$ is a compact metric space which is not an ANR.
By the classical theory of retracts (see, \emph{e.g.},~\cite[Chapter~IV, \S3]{borsuk}),
there exist a compact metric space~$Y$ and an embedding $X\hookrightarrow Y$ such that
$X$ is not a neighbourhood retract of~$Y$
(one may take $Y$ to be a suitable compact subset of the Hilbert cube).
Fix a compatible metric $d$ on~$Y$ and set
\[
Y_n:=\bigl\{y\in Y: d(y,X)\le 1/n\bigr\}\qquad (n\in\mathbb N).
\]
Then $Y_1\supseteq Y_2\supseteq\cdots\supseteq X$, each $Y_n$ is compact,
and $\bigcap_n Y_n=X$.
Since $X$ is not a neighbourhood retract of~$Y$,
there is no retraction from any $Y_n$ onto~$X$:
indeed, if $r_n\colon Y_n\to X$ were a retraction, then its restriction to
$W_n:=\{y\in Y:d(y,X)<1/n\}$ would be a retraction from an \emph{open} neighbourhood
$W_n$ of $X$ in $Y$, contradicting the choice of the embedding $X\hookrightarrow Y$.

Define $B:=C(Y)$, $I:=I_X$, and $I_n:=I_{Y_n}$.
Since $Y$ is compact metrisable, the Banach lattice $B=C(Y)$ is separable.
Lemma~\ref{lem:limit-ideal} gives $I=\overline{\bigcup_{n} I_n}$,
so $I$ is an inductive limit ideal in~$B$.
The canonical identifications are
\begin{equation}\label{eq:identifications}
B/I\cong C(X)\qquad\text{and}\qquad B/I_n\cong C(Y_n).
\end{equation}
Under these identifications, the canonical surjection
$\pi_n\colon B/I_n\to B/I$ corresponds to the restriction operator
$\rho\colon C(Y_n)\to C(X)$, $\rho(g)=g|_X$.

Assume, towards a contradiction, that $C(X)$ is semiprojective.
Let $\phi\colon C(X)\to B/I$ be the canonical isometric lattice isomorphism.
By Definition~\ref{def:semi} there exists $n\in\mathbb N$ such that for every $\varepsilon>0$
there is a lattice homomorphism $\psi_\varepsilon\colon C(X)\to B/I_n$
with $\|\psi_\varepsilon\|\le 1+\varepsilon$ and $\pi_n\circ \psi_\varepsilon=\phi$.
Via the identifications~\eqref{eq:identifications},
this yields for each $\varepsilon>0$ a lattice homomorphism
$T_\varepsilon\colon C(X)\to C(Y_n)$ satisfying
\begin{equation}\label{eq:section}
\rho\circ T_\varepsilon=\id_{C(X)}.
\end{equation}

Fix any such map $T:=T_\varepsilon$.
Apply Lemma~\ref{lem:weighted-composition} to~$T$: we obtain
an open set $U\subseteq Y_n$,
a~continuous weight $h=T\mathbf{1}\in C(Y_n)_+$,
and a continuous map $r\colon U\to X$ such that
$(Tf)(y)=h(y)\,f(r(y))$ for all $f\in C(X)$ and $y\in U$.

We claim that $r$ restricts to the identity on~$X$, so $r\colon U\to X$ is a retraction
from a neighbourhood of~$X$ in~$Y$ onto~$X$.
Indeed, by~\eqref{eq:section} we have $(T\mathbf{1})|_X=\mathbf{1}$, so
$h(x)=1$ for every $x\in X$.
In particular, $X\subseteq U=\{y\in Y_n:h(y)>0\}$.
For $x\in X$ and arbitrary $f\in C(X)$, \eqref{eq:section} gives
\[
f(x)=(Tf)(x)=h(x)\,f\bigl(r(x)\bigr)=f\bigl(r(x)\bigr).
\]
Separating points of~$X$ by continuous functions forces $r(x)=x$.
Set $W_n:=\{y\in Y: d(y,X)<1/n\}$.
Then $W_n$ is an \emph{open} neighbourhood of $X$ in $Y$ and $W_n\subseteq Y_n$.
Since $U$ is open in the subspace~$Y_n$, the intersection $U_0:=U\cap W_n$
is open in~$Y$, contains~$X$, and satisfies $U_0\subseteq U$.
Therefore $r|_{U_0}\colon U_0\to X$ is a continuous retraction from an open neighbourhood
of $X$ in $Y$ onto $X$, contradicting the choice of the embedding $X\hookrightarrow Y$.
\end{proof}

\section{Proof of Theorem~\ref{thm:B} and applications}\label{sec:thmB}

Having established a topological characterisation for $C(X)$-spaces,
it is natural to ask whether semiprojectivity and $1^+$-projectivity coincide
for general Banach lattices.
We recall the definition of the stronger property.

\begin{definition}[{\cite[Definition~10.1]{dpw}}]\label{def:projective}
A Banach lattice $P$ is \emph{$1^+$-projective} if for every Banach lattice~$X$,
every closed ideal $J\subseteq X$,
every lattice homomorphism $T\colon P\to X/J$, and every $\varepsilon>0$,
there exists a lattice homomorphism $\widehat T\colon P\to X$ with
$Q\circ \widehat T=T$ and $\|\widehat T\|\le (1+\varepsilon)\|T\|$,
where $Q\colon X\to X/J$ denotes the quotient map.
\end{definition}

\begin{proof}[Proof of Theorem~\ref{thm:B}]
If $\Gamma$ is countable, the $1^+$-projectivity of~$P$ follows from
de~Pagter and Wickstead~\cite[Theorem~11.8]{dpw}.

We now prove semiprojectivity for arbitrary $\Gamma$.
We begin with a general observation: every disjoint family of non-zero positive
elements in a separable Banach lattice is countable.
Indeed, let $E$ be a separable Banach lattice and let
$\{u_\gamma\}_{\gamma\in\Gamma}\subseteq E_+$
be a family of pairwise disjoint non-zero elements.
For each $m\in\mathbb N$, the set $\Gamma_m:=\{\gamma\in\Gamma: \|u_\gamma\|\ge 1/m\}$
is $1/m$-separated in~$E$:
if $\gamma\ne \gamma'$ lie in~$\Gamma_m$, then $u_\gamma\perp u_{\gamma'}$ gives
$|u_\gamma-u_{\gamma'}|=u_\gamma+u_{\gamma'}$, so
$\|u_\gamma-u_{\gamma'}\|=\|u_\gamma+u_{\gamma'}\|\ge \|u_\gamma\|\ge 1/m$.
Separability forces each $\Gamma_m$ to be countable,
and since $\Gamma=\bigcup_{m}\Gamma_m$, the claim follows.

Now let $B$ be a separable Banach lattice,
$I=\overline{\bigcup_n I_n}\subseteq B$ an inductive limit ideal,
and $\phi\colon P\to B/I$ a contractive lattice homomorphism.
For each $\gamma\in\Gamma$ let $\iota_\gamma\colon P_\gamma\to P$ denote the canonical coordinate
embedding, and fix a topological order unit $e_\gamma\in (P_\gamma)_+$.
Set
\[
v_\gamma:=\iota_\gamma(e_\gamma)\in P_+,
\qquad
u_\gamma:=\phi(v_\gamma)\in (B/I)_+.
\]
Since $v_\gamma\perp v_{\gamma'}$ in $P$ for $\gamma\neq\gamma'$ and $\phi$ is a lattice
homomorphism, the family $(u_\gamma)_{\gamma\in\Gamma}$ is pairwise disjoint in $B/I$:
\[
|u_\gamma|\wedge |u_{\gamma'}|
=|\phi(v_\gamma)|\wedge |\phi(v_{\gamma'})|
=\phi(v_\gamma\wedge v_{\gamma'})=0
\qquad(\gamma\neq\gamma').
\]
The quotient $B/I$, equipped with the quotient norm, is separable as a continuous image of~$B$,
so the set $J:=\{\gamma\in\Gamma: u_\gamma\ne 0\}$ is countable
by the observation above.

Now fix $\gamma\notin J$, so $u_\gamma=\phi(\iota_\gamma(e_\gamma))=0$.
Since $\ker\phi$ is a closed ideal in $P$ and contains $\iota_\gamma(e_\gamma)$,
it contains every $z\in P$ satisfying $|z|\le c\,\iota_\gamma(e_\gamma)$ for some $c>0$.
These are precisely the vectors of the form $\iota_\gamma(x)$ with
$|x|\le c\,e_\gamma$.
Hence
\[
\iota_\gamma(I(e_\gamma))\subseteq \ker\phi,
\qquad
I(e_\gamma):=\{x\in P_\gamma: |x|\le c\,e_\gamma \text{ for some } c>0\}.
\]
Because $e_\gamma$ is a topological order unit, $\overline{I(e_\gamma)}=P_\gamma$.
Since $\ker\phi$ is closed, we conclude that
\[
\iota_\gamma(P_\gamma)=\overline{\iota_\gamma(I(e_\gamma))}\subseteq \ker\phi.
\]
Consequently, $\phi(\iota_\gamma(x))=0$ for all $x\in P_\gamma$.

Let $P_J:=\bigl(\bigoplus_{\gamma\in J} P_\gamma\bigr)_{\ell_1(J)}$ and let $\pi_J\colon P\to P_J$
denote the coordinate projection (a contractive lattice homomorphism).
The previous paragraph shows that $\phi$ factors as
\[
P \xrightarrow{\;\pi_J\;} P_J \xrightarrow{\;\phi_0\;} B/I
\]
for a contractive lattice homomorphism~$\phi_0$.

Since $J$ is countable, the first paragraph of this proof gives that $P_J$ is
$1^+$-projective.
Let $Q\colon B\to B/I$ denote the quotient map.
For any $\varepsilon>0$, the $1^+$-projectivity of~$P_J$ furnishes
a lattice homomorphism $\widehat\phi_\varepsilon\colon P_J\to B$ satisfying
$Q\circ \widehat\phi_\varepsilon=\phi_0$ and
$\|\widehat\phi_\varepsilon\|\le 1+\varepsilon$.
Fix $n=1$ and let $Q_1\colon B\to B/I_1$ be the quotient map.
Define $\psi_\varepsilon:=Q_1\circ \widehat\phi_\varepsilon\circ \pi_J\colon P\to B/I_1$.
Then $\psi_\varepsilon$ is a lattice homomorphism with
$\|\psi_\varepsilon\|\le 1+\varepsilon$, and
\[
\pi_1\circ \psi_\varepsilon
=\pi_1\circ Q_1\circ \widehat\phi_\varepsilon\circ \pi_J
=Q\circ \widehat\phi_\varepsilon\circ \pi_J
=\phi_0\circ \pi_J
=\phi.
\]
Since the index $n=1$ does not depend on~$\varepsilon$,
$P$ is semiprojective.
\end{proof}

\begin{proof}[Proof of Corollary~\ref{cor:l1-Gamma-intro}]
The space $\ell_1(\Gamma)=\bigl(\bigoplus_{\gamma\in\Gamma}\mathbb R\bigr)_{\ell_1(\Gamma)}$
is a special case of Theorem~\ref{thm:B}
(each summand $\mathbb R$ is $1^+$-projective by~\cite[Corollary~10.6]{dpw},
with $e_\gamma=1$ as topological order unit).
If $\Gamma$ is countable, $1^+$-projectivity follows from Theorem~\ref{thm:B}.
If $\Gamma$ is uncountable, then $\ell_1(\Gamma)$ contains the uncountable
disjoint family $\{e_\gamma\}_{\gamma\in\Gamma}$;
by~\cite[Corollary~10.5]{dpw}, every $1^+$-projective Banach lattice has
only countably many pairwise disjoint non-zero positive elements,
so $\ell_1(\Gamma)$ is not $1^+$-projective.
\end{proof}

\begin{corollary}\label{cor:l1-sum-CK-anr}
Let\/ $\Gamma$ be a set and let\/ $(K_\gamma)_{\gamma\in\Gamma}$ be a family
of compact metric ANR's.
Then $\bigl(\bigoplus_{\gamma\in\Gamma} C(K_\gamma)\bigr)_{\ell_1(\Gamma)}$ is semiprojective.
If\/ $\Gamma$ is countable, then it is $1^+$-projective.
\end{corollary}

\begin{proof}
Each $C(K_\gamma)$ is $1^+$-projective by
Avil\'es--Mart\'{\i}nez-Cervantes--Rodr\'{\i}guez Abell\'an~\cite[Theorem~1.4]{avilaes2020},
with $\mathbf{1}_{K_\gamma}$ as topological order unit.
The result follows from Theorem~\ref{thm:B}.
\end{proof}

\begin{remark}\label{rem:l1-sum-why-special}
Theorem~\ref{thm:B} uses two ingredients that go beyond semiprojectivity
of the summands.
\emph{(1)~Order units detect support.}
The reduction to countably many coordinates relies on the existence, in each summand~$P_\gamma$,
of a topological order unit~$e_\gamma$ whose generated closed ideal is all of~$P_\gamma$.
If $\phi(\iota_\gamma(e_\gamma))=0$, then the closed-ideal property forces
$\phi$ to vanish on the entire coordinate $\iota_\gamma(P_\gamma)$.
\emph{(2)~Countable $\ell_1$-sums of $1^+$-projectives are $1^+$-projective.}
After the support reduction, one still needs a strong permanence theorem to lift the
resulting map from the countable subsum into~$B$ with uniform norm control,
provided by~\cite[Theorem~11.8]{dpw} for $1^+$-projective summands.
If one replaces the summands by merely \emph{semiprojective} Banach lattices,
neither ingredient is guaranteed:
semiprojectivity alone does not provide a canonical order unit,
and there is currently no analogue of~\cite[Theorem~11.8]{dpw} for semiprojective summands.
Concretely, semiprojectivity of~$A_\gamma$ yields lifts only at some stage $B/I_{n(\gamma)}$
(which may depend on~$\gamma$), and there is no reason for the indices $n(\gamma)$ to be
bounded when infinitely many coordinates are involved;
cf.\ Question~\ref{q:permanence}.
\end{remark}

\section{Concluding remarks and open questions}\label{sec:questions}

Theorem~\ref{thm:A} shows that semiprojectivity of $C(X)$ in the Banach-lattice category
captures the ANR property of~$X$ without the dimension-one constraint
that is present in the $C^*$-algebraic setting~\cite{sorensen-thiel}.
We close with several questions that arise naturally from this work.

\subsection*{$\ell$-open and $\ell$-closed Banach lattices}

Motivated by work of Oyetunbi and Tikuisis~\cite{oyetunbi-tikuisis}
on $\ell$-open and $\ell$-closed $C^*$-algebras,
one can introduce analogous notions in the Banach-lattice setting.

Fix Banach lattices $A,B$ and a closed ideal $I\subseteq B$.
A \emph{contractive} lattice homomorphism $\phi\colon A\to B/I$ is called
\emph{$1^+$-liftable} if for every $\varepsilon>0$ there exists a lattice homomorphism
$\widehat\phi_\varepsilon\colon A\to B$ such that $Q\circ \widehat\phi_\varepsilon=\phi$ and
$\|\widehat\phi_\varepsilon\|\le 1+\varepsilon$, where $Q\colon B\to B/I$ is the quotient map.

Let
\[
\Hom_1(A,B/I):=\bigl\{\phi\in\Hom(A,B/I): \|\phi\|\le 1\bigr\},
\]
endowed with the point-norm topology.
We say that $A$ is \emph{$\ell$-open} (respectively, \emph{$\ell$-closed}) if for every
separable Banach lattice $B$ and every closed ideal $I\subseteq B$, the set of
$1^+$-liftable elements of $\Hom_1(A,B/I)$ is open
(respectively, closed) in $\Hom_1(A,B/I)$.

Clearly, every $1^+$-projective Banach lattice is both $\ell$-open and $\ell$-closed.
In particular, if $X$ is a compact metric ANR, then $C(X)$ is $\ell$-open and $\ell$-closed.

In the \(C^*\)-algebraic setting, Oyetunbi and Tikuisis prove that a commutative unital
\(C^*\)-algebra is \(\ell\)-open if and only if it is semiprojective \cite{oyetunbi-tikuisis}.
Thus Question~\ref{q:l-open} asks whether the corresponding equivalence persists
in the Banach-lattice category. This problem appears more challenging than its $C^*-$algebraic counterpart: since lattice homomorphisms are not automatically contractive, techniques for transferring lifting properties must maintain strict norm control that could be easily destroyed under perturbation.

\begin{question}\label{q:l-open}
Do semiprojective Banach lattices have to be $\ell$-open?
Conversely, does $\ell$-openness imply semiprojectivity?
\end{question}

\subsection*{Further questions}

\begin{question}\label{q:separable-equivalence}
Let $A$ be a \emph{separable} Banach lattice.
Does semiprojectivity of $A$ imply that $A$ is $1^+$-projective?
Equivalently, do semiprojectivity and $1^+$-projectivity coincide for separable Banach lattices?
\end{question}

\begin{remark}
Corollary~\ref{cor:l1-Gamma-intro} separates semiprojectivity from $1^+$-projectivity using a nonseparable domain.
Question~\ref{q:separable-equivalence} asks whether nonseparability is essential for such a separation.
\end{remark}

\begin{question}\label{q:inductive-limit}
Is every separable Banach lattice an inductive limit of semiprojective Banach lattices?
\end{question}

This is the Banach-lattice counterpart of a question of
Blackadar~\cite[Section~4]{blackadar85} for $C^*$-algebras
(see also~\cite{thiel-inductive}).
In the commutative $C^*$-setting the answer is affirmative, since every compact metric space
is an inverse limit of ANR's,
and Thiel~\cite{thiel-inductive} established substantial closure properties
for the class of $C^*$-algebras that are inductive limits of semiprojective
$C^*$-algebras (for instance, closure under shape domination and homotopy equivalence).
It would be interesting to investigate whether analogous closure properties
hold in the Banach-lattice category,
and whether there is a shape-theoretic framework for Banach lattices
in which such questions can be studied systematically.

\begin{proposition}[Lifting through quotients of separable Banach lattices]\label{prop:sep-quotient-lift}
Let $A$ be a semiprojective Banach lattice, let $X$ be a separable Banach lattice,
let $J\subseteq X$ be a closed ideal, and let
$T\colon A\to X/J$ be a contractive lattice homomorphism.
Then there exists a lattice homomorphism $\widehat T\colon A\to X$ such that
$Q\circ \widehat T=T$, where $Q\colon X\to X/J$ is the quotient map.
\end{proposition}

\begin{proof}
Consider
\[
B:=\Bigl\{f\in C([0,1],X): f(0)=0
\text{ and } Q(f(t))=t\,Q(f(1)) \text{ for all } t\in[0,1]\Bigr\},
\]
equipped with the supremum norm and the pointwise order.
Then \(B\) is a closed sublattice of \(C([0,1],X)\): closedness follows because the
defining conditions are preserved under uniform limits, and if \(f\in B\), then
\[
Q(|f(t)|)=|Q(f(t))|=t\,|Q(f(1))|=t\,Q(|f(1)|)
\qquad (t\in[0,1]),
\]
so \(|f|\in B\).

Since $X$ is separable and $[0,1]$ is compact metric, $C([0,1],X)$ is separable,
hence so is $B$.

Define
\[
\Pi\colon B\to X/J,\qquad \Pi(f):=Q(f(1)).
\]
Then $\Pi$ is a contractive lattice homomorphism.
It is surjective: given $y\in X/J$, choose $x\in X$ with $Qx=y$ and set
$f_x(t):=tx$; then $f_x\in B$ and $\Pi(f_x)=y$.

Its kernel is
\[
I=\{f\in C([0,1],J): f(0)=0\}.
\]
For each $n\in\mathbb N$, let
\[
I_n:=\{f\in I: f|_{[0,1/n]}=0\}.
\]
Then each \(I_n\) is a closed ideal in \(B\), the sequence \((I_n)\) is increasing,
and \(I=\overline{\bigcup_n I_n}\). Indeed, closedness and monotonicity are clear. If
\(g\in I_n\) and \(f\in B\) satisfy \(|f|\le |g|\), then \(f(t)=0\) for \(t\in[0,1/n]\),
and since \(g(t)\in J\) and \(J\) is an ideal in \(X\), also \(f(t)\in J\) for every
\(t\in[0,1]\); hence \(f\in I_n\).
Indeed, if $f\in I$, choose $u_n\in C([0,1])$ with
$0\le u_n\le 1$, $u_n|_{[0,1/n]}=0$, and $u_n|_{[2/n,1]}=1$.
Then $u_nf\in I_n$ and
\[
\|f-u_nf\|_\infty\le \sup_{0\le t\le 2/n}\|f(t)\|\xrightarrow[n\to\infty]{}0,
\]
because $f(0)=0$.

The map $\Pi$ induces a lattice isomorphism
\[
\widetilde\Pi\colon B/I\to X/J.
\]
Moreover, $\widetilde\Pi$ is isometric: contractivity is clear, while for
$y\in X/J$ and $\eta>0$ one may choose $x\in X$ with $Qx=y$ and
$\|x\|\le \|y\|+\eta$; then $f_x(t)=tx$ belongs to $B$ and
\[
\|\widetilde\Pi^{-1}(y)\|\le \|f_x\|_\infty=\|x\|\le \|y\|+\eta.
\]
Letting $\eta\downarrow 0$ yields $\|\widetilde\Pi^{-1}(y)\|\le \|y\|$.

Now set
\[
\phi:=\widetilde\Pi^{-1}\circ T\colon A\to B/I.
\]
Since $A$ is semiprojective and $B$ is separable, there exist $n\in\mathbb N$
and, for instance with $\varepsilon=1$, a lattice homomorphism
\[
\psi\colon A\to B/I_n
\]
such that
\[
\pi_n\circ \psi=\phi,
\]
where $\pi_n\colon B/I_n\to B/I$ is the canonical quotient map.

Define
\[
E_n\colon B/I_n\to X,\qquad E_n([f]):=n\,f(1/n).
\]
This is well-defined because every element of $I_n$ vanishes at $1/n$.
It is a lattice homomorphism, since point evaluation and multiplication by a
positive scalar preserve lattice operations.
For $[f]\in B/I_n$ we have
\[
Q(E_n([f]))
=n\,Q(f(1/n))
=n\cdot \frac1n\,Q(f(1))
=\widetilde\Pi(\pi_n([f])).
\]
Hence
\[
Q\circ E_n\circ \psi
=\widetilde\Pi\circ \pi_n\circ \psi
=\widetilde\Pi\circ \phi
=T.
\]
Therefore $\widehat T:=E_n\circ \psi$ is the desired lift.
\end{proof}

\begin{remark}\label{rem:bounded-homogeneity}
By positive homogeneity, Definition~\ref{def:semi} and
Proposition~\ref{prop:sep-quotient-lift} remain valid with
`contractive' replaced by `bounded':
if $\phi$ has norm $M>0$, apply the contractive statement to $M^{-1}\phi$
and then multiply the resulting lift by $M$.
\end{remark}

\begin{corollary}\label{cor:ell-p-not-semi}
For every $1<p<\infty$, the Banach lattice $\ell_p$ is not semiprojective.
\end{corollary}

\begin{proof}
Let
\[
L:=\mathcal P_{\mathrm{fin}}(\mathbb N)\setminus\{\varnothing\},
\]
and let $FBL(L)$ denote the free Banach lattice generated by $L$ in the sense of
\cite{dpw}. Since $L$ is countable, $FBL(L)$ is separable by
\cite[Theorem~8.4]{dpw}.
By \cite[Lemma~2.2]{avilaes-c0}, there exists a surjective lattice homomorphism
\(
\Phi\colon FBL(L)\to c_0.
\)
Let
\(
J_p\colon \ell_p\to c_0
\) 
be the formal inclusion; this is a contractive lattice homomorphism.

Assume towards a contradiction that $\ell_p$ is semiprojective.
Let $Q\colon FBL(L)\to FBL(L)/\ker\Phi$ be the quotient map, and let
$\widetilde\Phi\colon FBL(L)/\ker\Phi\to c_0$
be the induced Banach lattice isomorphism.
Then $\widetilde\Phi^{-1}\circ J_p\colon \ell_p\to FBL(L)/\ker\Phi$
is a bounded lattice homomorphism.
By Remark~\ref{rem:bounded-homogeneity} and
Proposition~\ref{prop:sep-quotient-lift}, there exists a lattice homomorphism
$u\colon \ell_p\to FBL(L)$ such that
\[
Q\circ u=\widetilde\Phi^{-1}\circ J_p.
\]
Consequently,
\[
\Phi\circ u
=\widetilde\Phi\circ Q\circ u
=\widetilde\Phi\circ \widetilde\Phi^{-1}\circ J_p
=J_p.
\]
For each $n\in\mathbb N$, set
\[
f_n:=u(e_n)\in FBL(L)_+,
\]
where $(e_n)$ is the canonical basis of $\ell_p$.

For each $n\in\mathbb N$, define
\[
x_n^*:=\bigl(\chi_A(\{n\})\bigr)_{A\in L}\in[-1,1]^L,
\]
using the notation of \cite[Section~2]{avilaes-c0}.
Since
\[
\Phi(f_n)=J_p(e_n)=e_n,
\]
the definition of $\Phi$ gives
\[
f_n(x_n^*)=1
\qquad (n\in\mathbb N).
\]
Moreover, for every finite set $F\subseteq L$, one can choose
$n\notin \bigcup F$; then $x_n^*|_F=0$.
Therefore \cite[Lemma~2.3]{avilaes-c0} yields a subsequence
$(f_{n_k})_{k\ge 1}$ such that, for every $m\in\mathbb N$,
\[
\Bigl\|\sum_{k=1}^m f_{n_k}\Bigr\|\ge m-1.
\]
On the other hand,
\[
\Bigl\|\sum_{k=1}^m f_{n_k}\Bigr\|
=\Bigl\|u\Bigl(\sum_{k=1}^m e_{n_k}\Bigr)\Bigr\|
\le \|u\|\,\Bigl\|\sum_{k=1}^m e_{n_k}\Bigr\|_{\ell_p}
=\|u\|\,m^{1/p}.
\]
Since $p>1$, this is impossible for all $m$.
The contradiction shows that $\ell_p$ is not semiprojective.
\end{proof}

\begin{corollary}\label{cor:c0-not-semi}
The Banach lattice \(c_0\) is not semiprojective.
\end{corollary}

\begin{proof}
Let
\[
L:=\mathcal P_{\mathrm{fin}}(\mathbb N)\setminus\{\varnothing\},
\]
and let \(FBL(L)\) denote the free Banach lattice generated by \(L\).
Since \(L\) is countable, \(FBL(L)\) is separable by \cite[Theorem~8.4]{dpw}.
By \cite[Lemma~2.2]{avilaes-c0}, there exists a surjective lattice homomorphism
\[
\Phi\colon FBL(L)\to c_0.
\]

Assume towards a contradiction that \(c_0\) is semiprojective.
Let \(Q\colon FBL(L)\to FBL(L)/\ker\Phi\) be the quotient map, and let
\(\widetilde\Phi\colon FBL(L)/\ker\Phi\to c_0\)
be the induced Banach lattice isomorphism.
By Remark~\ref{rem:bounded-homogeneity} and
Proposition~\ref{prop:sep-quotient-lift}, there exists a lattice homomorphism
\(u\colon c_0\to FBL(L)\) such that
\[
Q\circ u=\widetilde\Phi^{-1}.
\]
Consequently,
\[
\Phi\circ u
=\widetilde\Phi\circ Q\circ u
=\id_{c_0}.
\]

For each \(n\in\mathbb N\), set \(f_n:=u(e_n)\in FBL(L)_+\), and define
\[
x_n^*:=\bigl(\chi_A(\{n\})\bigr)_{A\in L}\in[-1,1]^L.
\]
Since \(\Phi(f_n)=e_n\), the definition of \(\Phi\) gives
\[
f_n(x_n^*)=1
\qquad (n\in\mathbb N).
\]
Moreover, for every finite set \(F\subseteq L\), one can choose
\(n\notin \bigcup F\), and then \(x_n^*|_F=0\).
Therefore \cite[Lemma~2.3]{avilaes-c0}, applied with \(\varepsilon=1\),
yields a subsequence \((f_{n_k})_{k\ge 1}\) such that, for every \(m\in\mathbb N\),
\[
\Bigl\|\sum_{k=1}^m f_{n_k}\Bigr\|\ge m-1.
\]
On the other hand,
\[
\Bigl\|\sum_{k=1}^m f_{n_k}\Bigr\|
=\Bigl\|u\Bigl(\sum_{k=1}^m e_{n_k}\Bigr)\Bigr\|
\le \|u\|\,\Bigl\|\sum_{k=1}^m e_{n_k}\Bigr\|_{c_0}
=\|u\|,
\]
a contradiction.
\end{proof}

\begin{corollary}\label{cor:sequence-lattices-not-semi}
Let \(\Gamma\) be an infinite set.
Then \(c_0(\Gamma)\) is not semiprojective.
If \(1<p<\infty\), then \(\ell_p(\Gamma)\) is not semiprojective either.
\end{corollary}

\begin{proof}
Choose a countably infinite subset \(\Delta\subseteq\Gamma\).

For \(c_0(\Gamma)\), the coordinate embedding and coordinate projection furnish contractive
lattice homomorphisms
\[
i_0\colon c_0(\Delta)\to c_0(\Gamma), \qquad
r_0\colon c_0(\Gamma)\to c_0(\Delta),
\]
with \(r_0\circ i_0=\id_{c_0(\Delta)}\).
Hence semiprojectivity of \(c_0(\Gamma)\) would imply semiprojectivity of
\(c_0(\Delta)\cong c_0\) by Proposition~\ref{prop:retract}, contradicting
Corollary~\ref{cor:c0-not-semi}.

Now let \(1<p<\infty\). Again the coordinate embedding and coordinate projection yield
contractive lattice homomorphisms
\[
i_p\colon \ell_p(\Delta)\to \ell_p(\Gamma), \qquad
r_p\colon \ell_p(\Gamma)\to \ell_p(\Delta),
\]
with \(r_p\circ i_p=\id_{\ell_p(\Delta)}\).
Therefore semiprojectivity of \(\ell_p(\Gamma)\) would imply semiprojectivity of
\(\ell_p(\Delta)\cong \ell_p\) by Proposition~\ref{prop:retract}, contradicting
Corollary~\ref{cor:ell-p-not-semi}.
\end{proof}

\begin{proposition}\label{prop:finite-dim}
Every finite-dimensional Banach lattice is $1^+$-projective, and therefore semiprojective.
In particular, for every $n\in\mathbb N$ and every $1\le p\le \infty$,
the Banach lattice $\ell_p^n$ is semiprojective.
\end{proposition}

\begin{proof}
By \cite[Theorem~11.1]{dpw}, every finite-dimensional Banach lattice is projective
in the sense of \cite{dpw}, that is, $1^+$-projective in the terminology of
Definition~\ref{def:projective}.
Since $1^+$-projectivity implies semiprojectivity
(after lifting to $B$ one composes with any quotient map $Q_n\colon B\to B/I_n$),
the claim follows.
\end{proof}

\begin{proposition}\label{prop:hom-zero-obstruction}
Let \(E\) be a non-zero separable Banach lattice such that
\(\Hom(E,\mathbb R)=\{0\}\).
Then \(E\) is not semiprojective.
\end{proposition}

\begin{proof}
Assume towards a contradiction that \(E\) is semiprojective.
Let \(FBL[E]\) be the free Banach lattice generated by the underlying Banach space of \(E\),
and let \(\phi_E\colon E\to FBL[E]\) be the canonical linear isometry.
By \cite[Definition~2.2 and Theorem~2.4]{aviles-fblE}, the identity operator on \(E\)
extends to a surjective lattice homomorphism
\[
\beta_E\colon FBL[E]\to E
\]
satisfying \(\beta_E\circ \phi_E=\id_E\).
Set \(J:=\ker \beta_E\).

Since \(E\) is separable, so is \(FBL[E]\): if \(D\subseteq E\) is countable and dense, then
the closed sublattice generated by \(\phi_E(D)\) contains \(\phi_E(E)\), hence equals \(FBL[E]\).

Let \(Q\colon FBL[E]\to FBL[E]/J\) be the quotient map, and let
\(\widetilde\beta_E\colon FBL[E]/J\to E\) be the induced Banach lattice isomorphism.
By Remark~\ref{rem:bounded-homogeneity} and
Proposition~\ref{prop:sep-quotient-lift}, applied to the bounded lattice homomorphism
\(\widetilde\beta_E^{-1}\colon E\to FBL[E]/J\),
there exists a lattice homomorphism \(u\colon E\to FBL[E]\) such that
\[
Q\circ u=\widetilde\beta_E^{-1}.
\]
Consequently,
\[
\beta_E\circ u=\widetilde\beta_E\circ Q\circ u=\id_E.
\]

Now fix \(x^*\in E^*\).
By \cite[Theorem~2.4]{aviles-fblE}, there exists a lattice homomorphism
\(\widehat{x^*}\colon FBL[E]\to\mathbb R\) such that
\(\widehat{x^*}\circ \phi_E=x^*\).
Since \(u\) is a lattice homomorphism, so is \(\widehat{x^*}\circ u\colon E\to\mathbb R\).
Our assumption \(\Hom(E,\mathbb R)=\{0\}\) therefore yields
\[
\widehat{x^*}\circ u=0
\qquad (x^*\in E^*).
\]
Hence
\[
u(E)\subseteq \bigcap_{x^*\in E^*}\ker \widehat{x^*}.
\]

In the concrete realization of \(FBL[E]\) from \cite[Definition~2.2]{aviles-fblE},
elements of \(FBL[E]\) are positively homogeneous functions on \(E^*\), and
\(\widehat{x^*}(f)=f(x^*)\) for every \(f\in FBL[E]\).
Therefore
\[
\bigcap_{x^*\in E^*}\ker \widehat{x^*}=\{0\}.
\]
Thus \(u(E)=\{0\}\), contradicting \(\beta_E\circ u=\id_E\) and \(E\neq\{0\}\).
\end{proof}

\begin{corollary}\label{cor:Lp-unit-interval-not-semi}
For every \(1\le p<\infty\), the Banach lattice \(L_p([0,1])\) is not semiprojective.
\end{corollary}

\begin{proof}
The space \(L_p([0,1])\) is non-zero and separable.
By Proposition~\ref{prop:hom-zero-obstruction}, it suffices to know that
\(\Hom(L_p([0,1]),\mathbb R)=\{0\}\); this follows from \cite[Example 1(iii)]{dantas2022}.
\end{proof}

\begin{proposition}\label{prop:orlicz-heart-hom-zero}
Let \((\Omega,\Sigma,\mu)\) be an atomless measure space and let \(\Phi\) be an Orlicz
function. Denote by \(H^\Phi(\mu)\) the corresponding Orlicz heart,
\[
H^\Phi(\mu):=
\Bigl\{f\in L^0(\mu): \int_\Omega \Phi(|f|/\lambda)\,d\mu<\infty
\text{ for every }\lambda>0\Bigr\},
\]
equipped with the Luxemburg norm. Then
\[
\Hom\bigl(H^\Phi(\mu),\mathbb R\bigr)=\{0\}.
\]
In particular, if \(H^\Phi(\mu)\) is non-zero and separable, then \(H^\Phi(\mu)\) is not semiprojective.

Consequently, if \(L^\Phi(\mu)=H^\Phi(\mu)\), then
\[
\Hom\bigl(L^\Phi(\mu),\mathbb R\bigr)=\{0\}.
\]
If moreover \(L^\Phi(\mu)\) is non-zero and separable, then \(L^\Phi(\mu)\) is not semiprojective.
\end{proposition}

\begin{proof}
Suppose that \(\varphi\colon H^\Phi(\mu)\to\mathbb R\) is a non-zero lattice homomorphism.
Choose \(f\in H^\Phi(\mu)_+\) with \(\varphi(f)>0\), and let \(\varepsilon>0\).

Since \(f\in H^\Phi(\mu)\), the function
\[
g_\varepsilon:=\Phi(|f|/\varepsilon)
\]
belongs to \(L_1(\mu)\). Define a finite measure \(\nu_\varepsilon\) on \((\Omega,\Sigma)\) by
\[
\nu_\varepsilon(A):=\int_A g_\varepsilon\,d\mu
\qquad (A\in\Sigma).
\]
Because \(\mu\) is atomless and \(\nu_\varepsilon\ll\mu\), the measure \(\nu_\varepsilon\) is atomless.

Choose \(m\in\mathbb N\) such that \(\nu_\varepsilon(\Omega)/m\le 1\), and partition \(\Omega\)
into measurable sets \(A_1,\dots,A_m\) satisfying
\[
\nu_\varepsilon(A_k)=\frac{\nu_\varepsilon(\Omega)}{m}\le 1
\qquad (k=1,\dots,m).
\]
Set
\[
f_k:=f\mathbf 1_{A_k}\qquad (k=1,\dots,m).
\]
Then the \(f_k\)'s are pairwise disjoint, positive, and
\[
f=\sum_{k=1}^m f_k.
\]
Moreover,
\[
\int_\Omega \Phi(|f_k|/\varepsilon)\,d\mu
=\int_{A_k}\Phi(|f|/\varepsilon)\,d\mu
=\nu_\varepsilon(A_k)\le 1,
\]
hence \(\|f_k\|_\Phi\le \varepsilon\) for every \(k\).

Since \(\varphi\) is a lattice homomorphism, the numbers \(\varphi(f_k)\in\mathbb R_+\)
are pairwise disjoint. Therefore at most one of them is non-zero. On the other hand,
\[
\sum_{k=1}^m \varphi(f_k)=\varphi(f)>0,
\]
so there exists \(k_0\in\{1,\dots,m\}\) such that
\[
\varphi(f_{k_0})=\varphi(f).
\]
Consequently,
\[
0<\varphi(f)=\varphi(f_{k_0})
\le \|\varphi\|\,\|f_{k_0}\|_\Phi
\le \|\varphi\|\,\varepsilon.
\]
Since \(\varepsilon>0\) was arbitrary, this is impossible. Thus
\[
\Hom\bigl(H^\Phi(\mu),\mathbb R\bigr)=\{0\}.
\]

If \(H^\Phi(\mu)\) is non-zero and separable, Proposition~\ref{prop:hom-zero-obstruction}
shows that \(H^\Phi(\mu)\) is not semiprojective.

If \(L^\Phi(\mu)=H^\Phi(\mu)\), the asserted identity
\[
\Hom\bigl(L^\Phi(\mu),\mathbb R\bigr)=\{0\}
\]
is immediate. Under the additional assumptions that \(L^\Phi(\mu)\) is non-zero and separable,
the same proposition yields that \(L^\Phi(\mu)\) is not semiprojective.
\end{proof}

Proposition~\ref{prop:finite-dim} settles the finite-dimensional case,
while Corollary~\ref{cor:ell-p-not-semi} and
Corollary~\ref{cor:Lp-unit-interval-not-semi}
rule out the classical sequence lattices \(\ell_p\) \((1<p<\infty)\) and the function lattices
\(L_p([0,1])\) \((1\le p<\infty)\).
Proposition~\ref{prop:orlicz-heart-hom-zero} extends
\cite[Example~3.3(iii)]{dantas2022}. Recall also that \(H^\Phi(\mu)\) is the
order-continuous part of \(L^\Phi(\mu)\) by
\cite[Proposition~4.7]{gao-leung-xanthos};
thus the additional hypothesis \(L^\Phi(\mu)=H^\Phi(\mu)\) is precisely the
order-continuity of the Luxemburg norm.

\begin{question}\label{q:permanence}
Is the class of semiprojective Banach lattices closed under
suitable finite direct sums $($\emph{e.g.}, $\oplus_1$ or $\oplus_\infty$$)$?
More generally, which permanence properties
$($passage to complemented sublattices, quotients by closed ideals,
Fremlin tensor products$)$ does semiprojectivity enjoy?
As a basic permanence property, semiprojectivity passes to contractive retracts
$($Proposition~$\ref{prop:retract})$.
\end{question}

\begin{remark}
Even the case of the $\ell_1$-sum $A\oplus_1 B$ (or $\ell_\infty$-sum $A\oplus_\infty B$)
appears subtle:
a~lattice homomorphism $\phi\colon A\oplus_1 B\to E/I$ restricts to lattice homomorphisms
on $A$ and $B$ with disjoint ranges in the quotient, but semiprojective lifts of these restrictions
need not have disjoint ranges at a finite stage $E/I_n$.
\end{remark}

In the $C^*$-setting, Loring~\cite{loring} showed that finite direct sums of semiprojective
$C^*$-algebras are semiprojective (in the unital case),
and Blackadar~\cite{blackadar-survey} showed that matrix amplifications
preserve semiprojectivity.
It is natural to expect similar permanence in the Banach-lattice category,
but the available techniques differ substantially from the $C^*$-case.

On the negative side, some natural categorical operations destroy semiprojectivity.

\begin{remark}[Semiprojectivity is not preserved by pullbacks]\label{rem:pullback}
Let $K\subseteq[0,1]$ be the middle--third Cantor set and let
$\rho\colon C([0,1])\to C(K)$ be the restriction map.
Consider the pullback Banach lattice
\[
P:=\{(f,g)\in C([0,1])\oplus_\infty C([0,1]) : f|_K=g|_K\}.
\]
Then $C([0,1])$ is semiprojective $($Theorem~$\ref{thm:A})$, but $P$ is not semiprojective.
\end{remark}

\begin{proof}
Let \(q\colon [0,1]_0\sqcup[0,1]_1\to Z\) be the quotient map.
A continuous function on \(Z\) is equivalently a pair of continuous functions on the two copies
that agree on \(K\), hence \(P\) identifies isometrically (as a Banach lattice) with \(C(Z)\).

We claim that \(Z\) is not an ANR. Since \(Z\) is compact metric, it is enough to show that
\(Z\) fails to be locally contractible at each point of \(q(K)\); indeed, every metric ANR is
locally contractible (see, e.g., \cite{hanner-anr,borsuk}).

Fix \(x=q(t)\) with \(t\in K\), and let \(U\) be a neighbourhood of \(x\) in \(Z\).
Then \(q^{-1}(U)\) is an open subset of \([0,1]_0\sqcup[0,1]_1\) containing both copies of \(t\).
Hence there exists an open interval \(J\subseteq[0,1]\) containing \(t\) such that
\(J_0\cup J_1\subseteq q^{-1}(U)\), where \(J_i\subseteq[0,1]_i\) denotes the copy of \(J\).
Since \([0,1]\setminus K\) is a dense open union of complementary intervals, we may choose a
complementary interval \((a,b)\subseteq J\). Then
\[
L_{a,b}:=q([a,b]_0)\cup q([a,b]_1)\subseteq U,
\]
and \(L_{a,b}\) is a simple closed curve.

We next show that the inclusion \(i\colon L_{a,b}\hookrightarrow Z\) is not null-homotopic.
Identify \(S^1\) with the unit circle in \(\mathbb C\), and define continuous maps
\(F_0,F_1\colon[0,1]\to S^1\) by
\[
F_0(t)=F_1(t)=1 \quad (0\le t\le a),\qquad
F_0(t)=F_1(t)=-1 \quad (b\le t\le 1),
\]
and, for \(t\in[a,b]\),
\[
F_0(t):=\exp\!\Bigl(i\pi \frac{t-a}{b-a}\Bigr),\qquad
F_1(t):=\exp\!\Bigl(-i\pi \frac{t-a}{b-a}\Bigr).
\]
Because \((a,b)\cap K=\varnothing\), the maps \(F_0\) and \(F_1\) agree on \(K\); hence they
descend to a continuous map \(F\colon Z\to S^1\). On the circle \(L_{a,b}\), the restriction
\(F\circ i\) maps the arc \(q([a,b]_0)\) homeomorphically onto the upper semicircle and the arc
\(q([a,b]_1)\) homeomorphically onto the lower semicircle, with matching endpoints \(1\) and \(-1\).
Hence \(F\circ i\) has degree \(1\). Therefore \(F\circ i\) is not null-homotopic, and neither is \(i\).

Now let \(V\subseteq U\) be any neighbourhood of \(x\). As above, choose an open interval
\(J'\subseteq[0,1]\) containing \(t\) with \(J'_0\cup J'_1\subseteq q^{-1}(V)\), and then a
complementary interval \((a',b')\subseteq J'\). The corresponding simple closed curve
\[
L_{a',b'}:=q([a',b']_0)\cup q([a',b']_1)
\]
is contained in \(V\), and the inclusion \(L_{a',b'}\hookrightarrow U\) is not null-homotopic
by the previous paragraph. Since this map factors through \(V\hookrightarrow U\), the inclusion
\(V\hookrightarrow U\) cannot be null-homotopic.

Thus \(Z\) is not locally contractible at \(x\). Since \(x\in q(K)\) was arbitrary, \(Z\) is not
locally contractible, hence \(Z\) is not an ANR. By Theorem~\ref{thm:A}, \(C(Z)\cong P\) is not
semiprojective.
\end{proof}

\begin{remark}[Semiprojectivity is not preserved by quotients]\label{rem:quotients}
Semiprojectivity need not pass to quotients.
For instance, if $K\subseteq[0,1]$ is the Cantor set, then $C([0,1])$ is semiprojective
by Theorem~\ref{thm:A}, whereas $C(K)\cong C([0,1])/I_K$ is not semiprojective.
Consequently, semiprojectivity cannot be preserved by pushouts in general
(since quotients arise as special pushouts along the zero object).
\end{remark}

\begin{question}\label{q:extensions}
Let
\[
0\to J\to A\to B\to 0
\]
be a short exact sequence of Banach lattices.
If \(J\) and \(B\) are semiprojective, must \(A\) be semiprojective?
\end{question}

\begin{remark}
A positive answer to Question~\ref{q:extensions} would imply that semiprojectivity
is closed under finite \(\ell_1\)- and \(\ell_\infty\)-sums, by applying it to the split
exact sequences
\[
0\to J\to J\oplus_1 B\to B\to 0
\qquad\text{and}\qquad
0\to J\to J\oplus_\infty B\to B\to 0.
\]
Thus Question~\ref{q:extensions} is stronger than the finite direct-sum permanence
problem raised in Question~\ref{q:permanence}.

The pullback example of Remark~\ref{rem:pullback} shows that a negative answer is
plausible: if \(K\subseteq[0,1]\) is the Cantor set and
\[
P=\{(f,g)\in C([0,1])\oplus_\infty C([0,1]) : f|_K=g|_K\},
\]
then
\[
0\to I_K\to P\to C([0,1])\to 0
\]
is exact, \(P\) is not semiprojective, and \(C([0,1])\) is semiprojective.
What is presently unclear is whether the ideal \(I_K\) is semiprojective.
\end{remark}

\begin{remark}
In the $C^*$-algebraic setting semiprojectivity is formulated for separable
$C^*$-algebras, and the defining lifting property is unchanged if one restricts
the ambient algebra $B$ to be separable; see \cite[pp.~11--12]{blackadar-survey}.
We do not know whether the analogous reduction holds in the Banach-lattice category.

This suggests a density-sensitive variant of Definition~\ref{def:semi}.
For an infinite cardinal $\kappa$, one could call a Banach lattice $A$
\emph{$\kappa$-semiprojective} if the lifting property is required only for
Banach lattices $B$ with density character $<\kappa$.
If one wants the definition to reflect the ambient cardinal more faithfully,
one should probably also replace the sequential tower $(I_n)$ by an increasing
chain $(I_\alpha)_{\alpha<\lambda}$ of closed ideals indexed by an ordinal
$\lambda<\kappa$.
The present paper treats only the case $\kappa=\omega_1$, where
$\operatorname{dens}(B)<\omega_1$ means that $B$ is separable and countable
chains suffice.
\end{remark}

\subsection*{Acknowledgements}
We are grateful to Leonid Positselski (Prague) for drawing our attention to the
connection between semiprojectivity and finitely presentable objects
in the sense of Ad\'amek and Rosick\'y~\cite{adamek-rosicky};
see Remark~\ref{rem:finitely-presentable}.
The second-named author gratefully acknowledges support received from
NCN Sonata-Bis~13 (2023/50/E/ST1/00067).

\begin{thebibliography}{99}

\bibitem{adamek-rosicky}
J.~Ad\'amek and J.~Rosick\'y,
\emph{Locally presentable and accessible categories},
London Math. Soc. Lecture Note Ser., vol.~189,
Cambridge Univ. Press, Cambridge, 1994.

\bibitem{avilaes2020}
A.~Avil\'es, G.~Mart\'{\i}nez-Cervantes, and J.~D.~Rodr\'{\i}guez Abell\'an,
\emph{On projective {B}anach lattices of the form {$C(K)$} and {$FBL[E]$}},
J. Math. Anal. Appl. \textbf{489} (2020), no.~1, 124129.

\bibitem{avilaes-c0}
A.~Avil\'es, G.~Mart\'{\i}nez-Cervantes, and J.~D.~Rodr\'{\i}guez Abell\'an,
\emph{On the {B}anach lattice {$c_0$}},
Rev. Mat. Complut. \textbf{34} (2021), no.~1, 203--213.


\bibitem{aviles-lattice-projectivity}
A.~Avil\'es, G.~Mart\'{\i}nez-Cervantes, J.~D.~Rodr\'{\i}guez Abell\'an, and A.~Rueda Zoca,
\emph{Free Banach lattices generated by a lattice and projectivity},
Proc. Amer. Math. Soc. \textbf{150} (2022), no.~5, 2071--2082.

\bibitem{aviles-fblE}
A.~Avil\'es, J.~Rodr\'{\i}guez, and P.~Tradacete,
\emph{The free {B}anach lattice generated by a {B}anach space},
J. Funct. Anal. \textbf{274} (2018), no.~10, 2955--2977.

\bibitem{aviles-tradacete}
A.~Avil\'es and P.~Tradacete,
\emph{Amalgamation and injectivity in Banach lattices},
Int. Math. Res. Not. IMRN 2023, no.~2, 956--997.


\bibitem{blackadar85}
B.~Blackadar,
\emph{Shape theory for $C^*$-algebras},
Math.\ Scand. \textbf{56} (1985), 249--275.

\bibitem{blackadar-survey}
B.~Blackadar,
\emph{Semiprojectivity: a retrospective},
slides from a talk given in Copenhagen, May 2010.
Available at \href{http://bruceblackadar.com/Mathematics/copen10b.pdf}{\texttt{bruceblackadar.com/Mathematics/copen10b.pdf}}.

\bibitem{borsuk}
K.~Borsuk,
\emph{Theory of retracts},
Monografie Matematyczne, vol.~44,
Pa\'nstwowe Wydawnictwo Naukowe, Warsaw, 1967.

\bibitem{chigogidze-dranishnikov}
A.~Chigogidze and A.~N.~Dranishnikov,
\emph{Which compacta are noncommutative {AR}s?},
Topology Appl. \textbf{157} (2010), no.~4, 774--778.

\bibitem{dantas2022}
S.~Dantas, G.~Mart\'{\i}nez-Cervantes, J.~D.~Rodr\'{\i}guez Abell\'an,
and A.~Rueda Zoca,
\emph{Norm-attaining lattice homomorphisms},
Rev. Mat. Iberoam. \textbf{38} (2022), no.~3, 981--1002.

\bibitem{gao-leung-xanthos}
N.~Gao, D.~H.~Leung, and F.~Xanthos,
\emph{Duality for unbounded order convergence and applications},
Positivity \textbf{22} (2018), no.~3, 711--725.


\bibitem{hanner-anr}
O.~Hanner,
\emph{Some theorems on absolute neighborhood retracts},
Ark. Mat. \textbf{1} (1951), 389--408.

\bibitem{dpw}
B.~de~Pagter and A.~W.~Wickstead,
\emph{Free and projective {B}anach lattices},
Proc. Roy. Soc. Edinburgh Sect. A \textbf{145} (2015), no.~1, 105--143.

\bibitem{loring}
T.~A.~Loring,
\emph{Lifting solutions to perturbing problems in $C^*$-algebras},
Fields Institute Monographs, vol.~8,
Amer.\ Math.\ Soc., Providence, RI, 1997.

\bibitem{oyetunbi-tikuisis}
D.~Oyetunbi and A.~Tikuisis,
\emph{On $\ell$-open $C^*$-algebras and $\ell$-closed $C^*$-algebras},
preprint, \href{https://arxiv.org/abs/2310.15155}{\texttt{arXiv:2310.15155}}.

\bibitem{sorensen-thiel}
A.~P.~W.~S\o rensen and H.~Thiel,
\emph{A characterization of semiprojectivity for commutative $C^*$-algebras},
Proc. Lond. Math. Soc. (3) \textbf{105} (2012), no.~5, 1021--1046.

\bibitem{thiel-categories}
H.~Thiel,
\emph{Semiprojectivity and semiinjectivity in different categories},
Publ.\ Mat.\ Urug. \textbf{17} (2019), 263--274.

\bibitem{thiel-inductive}
H.~Thiel,
\emph{Inductive limits of semiprojective $C^*$-algebras},
Adv.\ Math. \textbf{347} (2019), 597--618.

\end{thebibliography}

\end{document}
